% ============================================================================
% CH — GROUNDING THE USER SIDE: fU-ProtoMF              [slug: fu-protomf]
% STATUS: design hardened through SC.5 (dc05, 2026-07-13); results BLOCKED on
% S0.7 fleet + SC.6+ runs. Lineage step 5. Datasets pinned: hm_1_month + ml-1m.
% DELTA CHAPTER (Matteo, 2026-07-14): fI (previous ch) is the full build —
% here concepts carry over and the chapter focuses on what is DIFFERENT,
% sometimes surprisingly so, and which NEW challenges arise. If fU were a
% precise mirror this chapter would be one sentence; it is not — the
% not-a-mirror surprises (mean normalization, B(t), empty histories,
% noise-for-crowding) ARE the content. Deliberately shorter than fI, and the
% structure reflects that: NO mirrors of fI's inventory / why-first /
% re-coordinatization sections (those transfer silently).
% HEADERS REWORKED 2026-07-14 (granularity): 6 -> 5 sections; the
% information-ordering punchline is now the CHAPTER INTRO (bonbon), not a
% section.
% Material map: ../outline.md#fu-protomf
% ============================================================================
\cleardoublepage
\chapter{Grounding the User Side: fU-ProtoMF (8--10)}   % page goal — scaffolding, DELETE before submission
\label{ch:fu}
\label{sec:fu-info-ordering}   % old opener section parked here (chapter intro)
\vault{2026-07-12_1713_information-ordering-uhost-fu-fi-same-info-different-side}

% CHAPTER INTRO (bonbon, unnumbered): the information-ordering punchline —
% info(U-ProtoMF) < info(fU) <= info(fI); the lineage holds information
% constant and moves which side it lives on. Frames the whole chapter as a
% delta against fI. JUSTIFICATION CALLBACK (Matteo, 2026-07-14): the intro
% justifies history-composition by calling back to fI's "Taste Is Revealed,
% Not Stated" section (\ref{sec:fi-why}) — users get composed from revealed
% behavior, not stated attributes, for exactly the reason argued there.

\section{The Model: History-Composed User Factors}
\label{sec:fu-model}
\vault{2026-07-12_1808_bt-implicit-item-intercept-shift-asymmetry-directives, 2026-07-12_1830_m5-noise-for-crowding-ids-division-of-labor-centering-option}
\vault{2026-07-14_1511_balog-scrutability-fu-linear-edit-affordance}
% 1511: Balog et al. (SIGIR 2019, weighted tag-set user model) = the closest
% TRANSPARENT ancestor of history composition — belongs in this model's
% positioning, not a generic RW pile. Scrutability = latent affordance of the
% linear sum (remove a purchase = subtract a row; enabled, not built; ID row
% = the non-scrutable residual). DIRECTIVE: "scrutability" and "Scrutiny
% Protocol" never appear near each other without the distinction drawn.
% DRIFT GUARD: do NOT re-teach composition mechanics (sum-of-rows, keystone
% reduction pattern, ablation-arm logic) — fI built all that; state "same
% construction, user side" in a sentence and spend the section on what
% DIFFERS: q_u = MEAN (not sum — why) of purchased items' feature-word
% embeddings (+ user-ID row) on the U-ProtoMF host; B(t) baseline mass;
% ids/noid arms; noise-for-crowding trade (M5'), division of labor
% (words = taste class, ID = idiosyncrasy).

\section{Users with Empty Histories: The Road to User-Attribute-Aware Prototypes}
\label{sec:fu-cold-users}
\label{sec:fufi-afu}   % afU outlook moved here from the merge chapter (2026-07-14)
\vault{2026-07-12_1653_dc05-zero-interaction-user-behavior-thesis-directive, 2026-06-16_1807_attribute-source-positioning}
% DIRECTIVE: exact and prominent treatment (own header kept for prominence;
% afU road promoted INTO the title — Matteo 2026-07-14: afU very likely WILL
% be built, so the title announces it rather than hiding it in an outlook).
% Title says "USER-attribute-aware" (Matteo 2026-07-14): everywhere else
% "attribute" means ITEM attributes — what afU adds is the user's own stated
% attributes; the noun is the destination capability (prototypes), the road
% itself is walked at factor level (cold-start fill).
% Content blocks:
%  1. Zero-interaction path: noise ID -> drop -> B(t) baseline fallback ->
%     fold-in recovery; graceful degradation of the aggregate.
%  2. THE ROAD (moved from the merge chapter — Matteo 2026-07-14, "part of
%     that line"; conscious revision of the committed lineage outline
%     2026-07-11_2351, step 7's presentation home): B(t) is a POPULATION
%     fallback — nothing user-specific fills a truly empty history; the
%     cold-start bell rings HERE. Candidate prime (afU): user attributes
%     enter last, weakest, only where nothing else exists — the placement
%     the taste argument demands (callback to the fI why-section). Expected
%     honest result stated up front. afU very likely WILL be built (Matteo
%     2026-07-14) — write as an announced road, not a hedge; if it ends up
%     unbuilt after all, the block degrades to Conclusion/Future Work.
%     Merged-model cold-user EMPIRICS, if any, stay in the merge chapter's
%     Results.

\section{Per-Purchase Attribution}
\label{sec:fu-attribution}
% The new explanation dimension (novelty axis: behavioral-evidence attribution):
% each recommendation decomposes over the user's own purchases / feature words.

\section{Results}
\label{sec:fu-results}
% [PENDING S0.7 + SC.6+] vs fleet; ids/noid arms; ml-1m results MUST carry the
% popularity-coupling mention (directive); two-arm spot-check.

\section{Explanations in Practice}
\label{sec:fu-explanations}
% DELTA ONLY (2026-07-14): system + host baseline live in the eval-framework
% chapter — here show only the new dimension (per-purchase regrouping /
% behavioral-evidence attribution), no host re-demo.
% Per-purchase regrouped explanation rendering; comparison with fI explanations.
